\documentclass[11pt]{article}
\usepackage[margin=1in]{geometry}
\usepackage{amsmath,amssymb,amsthm}
\usepackage{booktabs}
\usepackage{xcolor}
\usepackage{hyperref}

\newcommand{\farmfos}{\textsc{Farm-FOS}}
\newcommand{\gd}{\mathrm{gd}}        % growth-day map
\newcommand{\clip}{\operatorname{clip}}

\title{\farmfos: A Spatiotemporally-Weighted, Physics-Calibrated\\
Path Metric for Full-Season Farm Agents \\[4pt]
\large (math reference --- standalone, for internal review)}
\author{}
\date{}

\begin{document}
\maketitle

\noindent\textbf{Purpose of this note.} This is the math reference for our new
evaluation metric. It is deliberately self-contained so it compiles on its own.
Plain-English intuition lives in \texttt{farm\_fos\_explainer.md}; here we just
fix notation and write the equations precisely.

%==============================================================
\section{Setup and notation}

A scenario produces two action sequences over one soybean season.

\paragraph{Oracle (gold) workflow.}
\[
W^\star=\bigl(a^\star_1,\dots,a^\star_n\bigr),\qquad
a^\star_i=\bigl(\tau_i,\;R^\star_i,\;t^\star_i\bigr),
\]
where $\tau_i$ is the \emph{action type} (irrigate, apply\_fungicide,
spray\_pesticide, plant, harvest, \dots), $R^\star_i\subseteq\{0,\dots,63\}$ is
the \emph{target ridge set}, and $t^\star_i$ is the \emph{ideal time} of the
action (the oracle's timestamp).

\paragraph{Agent trace.}
\[
W=\bigl(b_1,\dots,b_m\bigr),\qquad b_k=\bigl(\tau_k,\;R_k,\;t_k\bigr),
\]
the analogous tuple for the agent's executed tool calls.

\paragraph{Physiological time.}
The crop responds to thermal time, not wall-clock time, so we measure all
offsets in \emph{growth-days} through a monotone map
\[
\gd:\;\text{sim-time}\;\longrightarrow\;\text{accumulated growing-degree-day index},
\]
i.e.\ $\gd(t)$ is the cumulative GDD (base $10^\circ$C, capped $30^\circ$C) up to
time $t$ for that scenario's weather trace. Writing offsets as
$\gd(t_k)-\gd(t^\star_i)$ makes ``three days late'' mean the same physiological
thing in a cold spring as in a hot August.

%==============================================================
\section{The three local ingredients}

For each oracle action $a^\star_i$ we score the agent on \emph{space}, \emph{time},
and \emph{physical importance}.

\subsection{Spatial overlap (where)}
\[
o\bigl(R,R^\star\bigr)=\frac{\lvert R\cap R^\star\rvert}{\lvert R^\star\rvert}\in[0,1].
\]
This is the fraction of the intended ridges the agent actually treated.
Whole-field / non-spatial operations take $R^\star=\{0,\dots,63\}$, so $o=1$ when
applied field-wide.

\subsection{Temporal kernel (when)}
The professor's ``local-maximum neighbourhood'': full credit at the ideal time,
decaying as the action drifts early or late. Primary form is the triangular
(``linear-rule'') kernel
\[
k_\sigma(\Delta)=\max\!\Bigl(0,\;1-\tfrac{\lvert\Delta\rvert}{\sigma}\Bigr),
\qquad \Delta=\gd(t_k)-\gd(t^\star_i),
\]
with action-specific tolerance (half-width) $\sigma$. A smooth Gaussian variant
\(
k_\sigma(\Delta)=\exp\!\bigl(-\Delta^2/2\sigma^2\bigr)
\)
is reported as an ablation.

\paragraph{Harvest is asymmetric.} Harvesting a few days \emph{late} (within the
post-R8 grace window) costs almost nothing, but harvesting \emph{before} R8 is
severe, and very late triggers shatter loss. We therefore use
\[
k^{\mathrm{harv}}(\Delta)=
\begin{cases}
1, & 0\le \Delta\le G \quad(\text{grace, }G\approx7\text{ d}),\\[2pt]
\max\!\bigl(0,\,1-\tfrac{\Delta-G}{\sigma_{\mathrm{late}}}\bigr), & \Delta>G,\\[4pt]
\max\!\bigl(0,\,1-\tfrac{\lvert\Delta\rvert}{\sigma_{\mathrm{early}}}\bigr), & \Delta<0,
\end{cases}
\qquad \sigma_{\mathrm{early}}\ll\sigma_{\mathrm{late}} .
\]

\subsection{Physical importance (how much it matters)}
$w_i\ge 0$ is the \emph{marginal yield fraction} attributable to action $i$ in
that scenario. It is read from the growth model (Section~\ref{sec:calib}), so a
pod-fill irrigation in a drought scenario carries large $w_i$ while a routine
check carries small $w_i$.

%==============================================================
\section{Matching agent actions to oracle actions}

Each oracle action is matched to the single agent action of the same type that
best aligns in space and time:
\[
\pi(i)=\arg\max_{\,k:\;\tau_k=\tau_i}\;
o\bigl(R_k,R^\star_i\bigr)\,k_{\sigma_i}\!\bigl(\gd(t_k)-\gd(t^\star_i)\bigr).
\]
Matching is greedy in descending credit so each agent action is consumed at most
once; if no same-type action exists (or the best alignment is $0$) the oracle
action is \emph{omitted} and earns zero credit. Write
\[
o_i=o\bigl(R_{\pi(i)},R^\star_i\bigr),\qquad
\Delta_i=\gd(t_{\pi(i)})-\gd(t^\star_i).
\]

%==============================================================
\section{The metric}

\begin{equation}
\boxed{\;
\farmfos\bigl(W,W^\star\bigr)=
\clip_{[0,1]}\!\left(
\frac{\displaystyle\sum_{i=1}^{n} w_i\,o_i\,k_{\sigma_i}(\Delta_i)
\;-\;\lambda\!\!\sum_{k\in U} h_k}
{\displaystyle\sum_{i=1}^{n} w_i}
\right)\;}
\label{eq:farmfos}
\end{equation}

\noindent where $U$ is the set of agent actions matched to no oracle action and
$h_k\ge 0$ is the physical \emph{harm} of such a spurious action (e.g.\ an
off-window spray that wastes a residual window, or needless irrigation that
leaches nutrients); most spurious actions are harmless, $h_k\approx0$. The scalar
$\lambda$ scales the harm penalty.

\paragraph{Reading the equation.} Each oracle action contributes its physical
weight $w_i$ scaled by how well the agent hit it in \emph{space} ($o_i$) and
\emph{time} ($k_{\sigma_i}$). The denominator normalises per scenario, so
$\farmfos\in[0,1]$ and the scenario's own stress profile decides which actions
dominate. Doing the right thing at the right place and the right moment
$\Rightarrow$ score $\to 1$; right action but mistimed around a high-$w$ window
$\Rightarrow$ proportional loss; missing a high-$w$ action $\Rightarrow$ large
drop.

%==============================================================
\section{Physics calibration of $w_i$ and $\sigma_i$}
\label{sec:calib}

\paragraph{Step 1 --- a-priori agronomic defaults (the clean rule).}
$w_i=\omega(\tau_i,s_i)$ and $\sigma_i=\varsigma(\tau_i,s_i)$ are read from a
small table indexed by action type and the growth stage $s_i$ at $t^\star_i$.
The timing windows are \emph{published crop-stage agronomy} (critical weed-free
period VE--V3; R5 as the most water-sensitive stage; R3--R6 disease window;
$13$--$15\%$ harvest-moisture window) --- the same a Harbin soybean operation
would use, independent of our simulator --- and the engine coefficients below
are shown to \emph{agree with}, not \emph{define}, them. Crucially these
parameters are \textbf{frozen before any yield correlation is computed}:

\begin{center}\small
\begin{tabular}{llll}
\toprule
Action & Peak stage & $\sigma$ (growth-days) & weight basis (engine) \\
\midrule
irrigate     & R5 pod-fill ($\pm$R1--R3)     & $\sim$7 (R5), $\sim$14 (veg) & water-stress $0.14$--$0.18$ VWC $\to$ RUE$^{\text{ss}}$ \\
fungicide    & R3--R6                        & $\sim$7  & LAI loss $0.055/\text{d}$ if disease$>0.25$; $0.45$ weight \\
insecticide  & V4$^+$--R5                    & $\sim$7  & $0.32$ biotic weight; suitability $=1$ V3--R5 \\
herbicide    & VE--V3                        & $\sim$3  & $0.28$ biotic weight; weeds escalate $0.049/\text{d}$ \\
fertilize    & V3--V4$^+$,\,R3--R5           & $\sim$14--21 & nutrient index $0.35$--$0.85$; slow $0.0008/\text{d}$ decay \\
plant        & VE soil-window               & $\sim$5  & stand fraction multiplies whole-season LAI \\
harvest      & R8$+7$..$14$ (asym.)          & grace $7$ $+\,0.3\%/\text{d}$ shatter & drydown $1\%/\text{d}$ + moisture machine-loss \\
\bottomrule
\end{tabular}
\end{center}

\paragraph{Step 2 --- counterfactual \emph{corroboration} (not fitting).}
As an after-the-fact sanity check on a handful of decisive actions, we compare
the a-priori parameters against the engine's own sensitivity. Define the
counterfactual yield drop from removing action $i$ from the oracle and replaying
to maturity,
\[
w_i^{\mathrm{cf}}=\frac{Y^\star-Y^{\star\setminus i}}{Y^\star},
\]
and sweep the action's day to trace the \emph{yield-vs-offset} curve. We show its
peak and half-width are \emph{consistent with} our a-priori $t^\star_i,\sigma_i$.
We \textbf{do not set $w_i:=w_i^{\mathrm{cf}}$}: equating the metric weights with
the model's own yield deltas would make \farmfos{} a reconstruction of $Y$ and
its correlation with $Y$ tautological. These curves appear in the appendix as
corroboration that the agronomy priors point the right way --- not as the source
of the weights.

\begin{remark}[Why this is not circular]
All four metrics in Section~\ref{sec:baselines} are functions of the same
$(W,W^\star)$ trace pair. \farmfos{} differs only in \emph{using} the timing,
spatial, and importance information that the baselines discard, with parameters
fixed a priori from agronomy. Any correlation gain therefore reflects which
information is yield-relevant, not access to the yield being predicted. We target
a strong-but-imperfect correlation (Spearman $\approx 0.6$--$0.85$); a
near-perfect value would indicate accidental reconstruction and is treated as a
diagnostic to investigate, not a result to report.
\end{remark}

%==============================================================
\section{Baselines as degenerate cases}
\label{sec:baselines}

The metrics we compare against are special cases of Eq.~\eqref{eq:farmfos} that
throw away one or more ingredients:

\begin{center}\small
\begin{tabular}{lcccc}
\toprule
Metric & spatial $o_i$ & temporal $k$ & weights $w_i$ & what it keeps \\
\midrule
BFCL success rate     & $1$ & $1$ & uniform & did the right tools get called \\
CORE / path-corr.\    & $1$ & $1$ & uniform & type \emph{order} (Levenshtein) \\
vanilla KTC           & $1$ & $1$ & uniform & type \emph{order} (Kendall $\tau$) \\
\textbf{\farmfos}     & $o_i$ & $k_{\sigma_i}(\Delta_i)$ & physics & order \emph{+ where + when + how much} \\
\bottomrule
\end{tabular}
\end{center}

\noindent Concretely, with $o_i\equiv1$, $k\equiv1$, $w_i$ uniform and
$\lambda=0$, Eq.~\eqref{eq:farmfos} collapses to
$\tfrac1n\sum_i\mathbb{1}[\text{matched}_i]$, i.e.\ a BFCL-style success rate.
CORE and KTC additionally score the \emph{order} of the matched types but remain
blind to the clock and to physical importance. \farmfos{} restores continuous
timing ($k$), spatial targeting ($o$), and physics weighting ($w$) --- the three
things that determine yield --- which is the empirical claim we test by
correlating each metric against measured season yield.

%==============================================================
\section{Evaluation protocol (what the scatter shows)}

For each of the $N$ runs we compute the four metrics above and the measured
outcome (preserved-yield ratio, equivalently $1-\text{crop\_loss}$). We plot
metric (y) vs.\ yield (x) and report Pearson $r$ and Spearman $\rho$. We test
that \farmfos's correlation exceeds each baseline's with a Steiger/Williams test
for dependent correlations. The hypothesis is
\[
\bigl|\rho_{\farmfos,\,Y}\bigr|\;>\;
\max\bigl(|\rho_{\mathrm{BFCL},Y}|,\,|\rho_{\mathrm{CORE},Y}|,\,|\rho_{\mathrm{KTC},Y}|\bigr),
\]
with the baselines already observed to be near zero on our $90$-run L3 set.

\paragraph{Two correlation regimes.} We report this in two modes that answer
different questions. \emph{(i) Across scenarios, fixed agent} (the $90$-run L3
set): does the metric track per-scenario performance of one agent? \emph{(ii)
Across agents, fixed scenario} (the $10$-controller $\times\ 3$-model set from
prior runs): on the same field, does a better agent earn both higher \farmfos{}
and higher yield? Regime (ii) is what shows the metric measures \emph{agent
skill}, not merely scenario difficulty, and pre-empts the confound that yield
and metric both track how hard a scenario is.

%==============================================================
\section{Limitations (stated up front)}

(1) \textbf{Simulated outcome.} $Y$ is produced by the same process model the
metric is calibrated from; this is a benchmark-wide property (all yields here are
simulated), and real-field validation is future work. (2) \textbf{Domain
calibration.} The numeric weights/windows are specific to Harbin soybean under
this engine; the transferable contribution is the \emph{recipe} (calibrate the
evaluation metric from the domain process model / known agronomy), not the
constants. (3) \textbf{Model-belief inheritance.} The metric measures fidelity
to the \emph{modeled} agronomy, which is also what defines the oracle and the
yield, so it is internally self-consistent by construction; we do not claim it
recovers ground-truth field agronomy. We argue a trace metric remains useful
even where $Y$ is computable, because it gives per-action diagnostic credit and
does not require a full-season rollout the evaluator may be unable to run.

\end{document}
